\documentclass[11pt]{article}
\usepackage[margin=1in]{geometry}
\usepackage{amsmath,amssymb}
\usepackage{booktabs}
\usepackage{graphicx}
\usepackage[hidelinks]{hyperref}
\usepackage{times}

\title{\textbf{NeuroCardia}\\[4pt]
\large An Open-Source, On-Device Phase-Domain Diagnostic Engine for
Single-Lead ECG\\[4pt]
\normalsize White Paper --- Technical and Regulatory Overview\\[2pt]
\small Version 0.2 --- Revision}
\author{Americo Simoes \\ \texttt{amexsimoes@gmail.com}}
\date{October 14, 2026}

\begin{document}
\maketitle

\begin{center}
\small
Copyright \textcopyright{} 2026 Americo Simoes. All rights reserved.\\
This white paper is released under the Creative Commons Attribution 4.0
International License (CC-BY-4.0). The software described in this
document is released under the GNU General Public License, version 3 or
later (GPL-3.0-or-later).
\end{center}

\vspace{1em}
\noindent\fbox{\parbox{\dimexpr\linewidth-2\fboxsep-2\fboxrule}{%
\textbf{Revision history.} This is revision 0.2 of the NeuroCardia
White Paper. Revision 0.1 (October 14, 2026) reported the pilot results
of the Kuramoto residual (AUC 0.89 on a paired corpus of nine sinus and
nine AF recordings) and the trajectory-correlation outcome of the
five-arm sensitivity analysis. Revision 0.2 adds: the residual AUC for
all five experimental arms (two of which do not exhibit the residual);
the orthogonality result (the residual retains AUC 0.716 after
regressing out $\mathrm{sd1}/\mathrm{sd2}$ and the coefficient of
variation); the finding that trajectory robustness and residual
robustness are disjoint; and a scope condition on the residual's
frequency-distribution dependence. The regulatory strategy is
unchanged. The list of claims requiring additional validation now
includes residual orthogonality to standard HRV.}}

\vspace{1em}

\begin{abstract}
NeuroCardia is a free, open-source Android application for offline,
on-device analysis of single-lead electrocardiogram (ECG) recordings.
It requires no network connection, no user accounts, and no telemetry;
every computation runs on the handset, and the application functions
identically in airplane mode. The application computes conventional
heart rate variability (HRV) metrics, a four-axis signal quality score,
and a phase-domain Cardiac Synchronization Index (CSI) derived from a
Kuramoto coupled-oscillator network driven by the RR tachogram. Its DSP
and I/O layers are packaged as standalone Kotlin/JVM modules that can be
reused in any JVM project.

This document specifies the computational model, the diagnostic
metrics, the validation strategy against public reference datasets, and
a preliminary regulatory pathway toward FDA clearance as a Class II
medical device under 510(k). It is explicit about which claims are
literature-supported, which are novel and awaiting validation, and
which are research hypotheses. The CSI is presented as a research
feature, not a clinical diagnostic; the application's primary
contribution is the deployment envelope --- open-source, offline,
on-device ECG analysis on commodity Android hardware. The source code
is released under GPL-3.0-or-later.
\end{abstract}

\tableofcontents

\section{Introduction}

\subsection{Motivation}

Heart rate variability analysis has been clinically accepted for
decades as a window into autonomic function~\cite{taskforce1996}. The
algorithms are mature and well-understood. What is not widely available
is a \emph{deployment envelope} in which those algorithms can be run
without a workstation, without a network connection, and without
sending patient data to a third party.

The dominant open-source tools for ECG and HRV analysis ---
predominantly Python-based research stacks --- assume a workstation, a
network connection, and a user comfortable with a command line. Field
clinicians, rural health workers, air-gapped research environments, and
privacy-sensitive deployments are left without validated, offline
alternatives. Commercial Holter analysis software is closed-source and
typically tied to a specific vendor's hardware. The gap is not in the
algorithms; it is in the delivery.

NeuroCardia addresses that gap with four design commitments.

\begin{enumerate}
\item \textbf{On-device only.} No cloud, no accounts, no telemetry. The
application functions identically in airplane mode. Every computation,
from R-peak detection to diagnostic report generation, runs on the
handset.
\item \textbf{Open source.} GPL-3.0-or-later. Every algorithm is
auditable, reproducible, and forkable. There is no proprietary binary
blob and no vendor lock-in.
\item \textbf{Reusable modules.} The DSP and I/O layers are packaged as
pure-Kotlin/JVM libraries with no Android dependency. They can be
consumed by any JVM project and tested in under a second, independent
of the Android application that uses them.
\item \textbf{Standard metrics first, novel metrics clearly labeled.}
Time-domain and frequency-domain HRV metrics are computed using
standard formulations and are the application's validated outputs. The
phase-domain Cardiac Synchronization Index is presented as a research
feature whose record-level diagnostic behavior is documented in the
companion residual paper and whose clinical utility is not yet
established.
\end{enumerate}

The primary contribution of the project is not a new algorithm. It is
the combination of standard algorithms with a deployment envelope that
does not require network access, third-party services, or a specific
hardware vendor.

\subsection{Scope of this document}

This white paper specifies: the computational model
(Section~\ref{sec:model}), the diagnostic metrics and their intended
clinical meaning (Section~\ref{sec:metrics}), the validation strategy
against public datasets (Section~\ref{sec:validation}), the regulatory
pathway toward FDA clearance (Section~\ref{sec:regulatory}), the
limitations and open questions (Section~\ref{sec:limitations}), and the
development roadmap (Section~\ref{sec:roadmap}).

It is written for three audiences: clinical collaborators evaluating
the diagnostic claims, regulatory reviewers assessing the device
pathway, and software engineers who may contribute to the
implementation.

\subsection{What this document is not}

This is not a clinical validation study. It does not report prospective
trial data, and it makes no claim of diagnostic accuracy against a gold
standard. It is a design and rationale document. Validation is
described as a strategy (Section~\ref{sec:validation}), not as a
result.

\subsection{What this application is not}

NeuroCardia is not a replacement for standard ECG interpretation, and
it is not a diagnostic device. Three specific exclusions are stated
here because they shape every subsequent claim.

\begin{itemize}
\item \textbf{It does not diagnose.} It computes descriptive metrics
from an ECG recording. It does not produce a diagnosis, does not
recommend treatment, and does not triage patients. Every report
produced by the application includes a disclaimer that the output is
informational and requires review by a qualified clinician.
\item \textbf{The CSI is a research feature.} The Cardiac
Synchronization Index and its derived early warning signals are not
validated diagnostics. They are documented in this paper and in the
companion residual paper as record-level descriptors whose clinical
significance is a subject of ongoing investigation. They are gated
behind a settings flag and labeled as not-for-clinical-use.
\item \textbf{It is not a real-time monitor.} The application analyzes
completed ECG recordings loaded from files. It does not stream from a
device, does not monitor continuously, and does not alarm. The
phase-domain metric in particular requires the full recording to
stabilize and is not a short-window detector.
\end{itemize}

These exclusions are deliberate and are not limitations of the current
version. They define the intended use. A device that made any of these
three claims would be a different device with a different regulatory
pathway and a different validation burden.

\section{Background}

\subsection{Heart rate variability}

HRV analysis quantifies beat-to-beat variations in the RR interval
series. Standard metrics include time-domain measures such as the root
mean square of successive differences (RMSSD), the standard deviation
of normal-to-normal intervals (SDNN), and the proportion of successive
differences greater than 50 ms (pNN50); and frequency-domain measures
such as low-frequency power (LF, 0.04--0.15 Hz), high-frequency power
(HF, 0.15--0.40 Hz), and the LF/HF ratio~\cite{taskforce1996}.

These metrics are descriptive: they characterize the state of the
autonomic nervous system at a moment in time. They do not, by
themselves, describe how that state is evolving, and they do not
surface the dynamical signatures of impending transitions.

\subsection{Cardiac rhythm as a coupled oscillator system}

The sinoatrial node is a population of coupled pacemaker cells whose
collective firing produces the cardiac rhythm. This is mathematically a
system of coupled phase oscillators, and it is naturally described by
the Kuramoto model~\cite{kuramoto1975,acebron2005}:
\begin{equation}
\frac{d\theta_i}{dt} = \omega_i + \frac{K}{N}\sum_{j=1}^{N}
\sin(\theta_j - \theta_i),
\label{eq:kuramoto}
\end{equation}
where $\theta_i$ is the phase of oscillator $i$, $\omega_i$ its natural
frequency, $K$ the coupling strength, and $N$ the number of
oscillators. The degree of synchronization is captured by the order
parameter:
\begin{equation}
r = \left| \frac{1}{N} \sum_{j=1}^{N} e^{i\theta_j} \right|,
\label{eq:orderparam}
\end{equation}
with $r \to 1$ indicating full phase locking and $r \to 0$ indicating
incoherence.

\subsection{Prior work on phase analysis of cardiac rhythm}

Phase synchronization indices applied to ECG have been shown to
outperform correlation, coherence, and mutual information for detecting
rotors during atrial fibrillation. Early warning signals derived from
Kuramoto-type models have been demonstrated to detect impending
transitions to fibrillation in computational studies. What is missing
is an integrated, open-source, on-device implementation.

\section{Computational Model}\label{sec:model}

\subsection{RR-to-phase mapping}

Given an RR interval series $\{R_k\}_{k=1}^{M}$ (in seconds), we
construct a phase trajectory by cumulative integration:
\begin{equation}
\phi_k = 2\pi \sum_{j=1}^{k} \frac{\overline{R}}{R_j} \pmod{2\pi},
\label{eq:rrphase}
\end{equation}
where $\overline{R}$ is the mean RR interval. Each beat advances the
phase by $2\pi \overline{R}/R_j$, so a beat at the mean rate advances
exactly $2\pi$. This maps the irregular cardiac rhythm onto a unit
circle and preserves the relative timing of each beat.

\subsection{Network construction}

We instantiate a Kuramoto network of $N = 16$ oscillators. Natural
frequencies are assigned deterministically from the non-trivial
Riemann zeros $\gamma_i$, then normalized to zero mean and unit
variance:
\begin{equation}
\omega_i = \frac{\gamma_i - \bar\gamma}{\sigma_\gamma},
\label{eq:freqnorm}
\end{equation}
where $\bar\gamma$ and $\sigma_\gamma$ are the mean and standard
deviation of the first 16 Riemann zeros. The normalization is required
for numerical tractability: the reference parameterization places the
carrier at $\Omega_T = 587{,}000$ rad/s, five orders of magnitude above
any frequency the Euler integrator can resolve at a practical time
step.

The frequency distribution is a specific modeling choice. Section~5.4
reports the sensitivity of the diagnostic to this choice. The full
sensitivity analysis is in the companion document~\cite{simoes2026sensitivity}.

The coupling strength is $K = 2$. This value is not arbitrary: as
shown in Section~5.4, the CSI at $K = 2$ carries incremental diagnostic
information over conventional HRV, while at higher coupling strengths
the network saturates and the metric becomes a monotone function of RR
variability.

\subsection{Phase reset from beats}

At each detected R-peak, we reset the phase of a dedicated heart
oscillator (index 0) to the cumulative phase $\phi_k$ from
Equation~\ref{eq:rrphase}. The remaining $N-1$ oscillators evolve under
Equation~\ref{eq:kuramoto}. This reset is what couples the model to
the actual cardiac signal.

\subsection{Coupling topologies}

Four coupling topologies are supported, selectable from a settings
menu: all-to-all (default), nearest-neighbor, small-world, and star.
The topology is a modeling choice, not a diagnostic one.

\subsection{Order parameter trajectory}

We integrate the network with a fixed step $\Delta t = 0.05$ s and
record the order parameter $r$ at each step. The result is a time
series $r(t)$ that we term the Cardiac Synchronization Index (CSI).

\section{Diagnostic Metrics}\label{sec:metrics}

\subsection{Cardiac synchronization index}

The CSI $r(t)$ is the primary novel output. It is bounded in $[0,1]$ by
construction. In the production configuration ($K = 2$), the mean value
across the reference corpora falls in the range 0.60--0.70; at higher
coupling strengths the metric saturates and sits at 0.98--0.99 for all
rhythms. The interpretation of specific CSI values is therefore not
fixed across parameterizations, and the metric should be read as a
relative quantity within a fixed parameterization, not as an absolute
one.

The trajectory of $r(t)$ --- its derivative, variance, and higher
moments --- is as diagnostic as its instantaneous value. The app
reports both the instantaneous value and the trajectory statistics.

The CSI is a parameterized quantity. Its value depends on the frequency
distribution of the oscillator network, the coupling strength $K$, and
the number of oscillators $N$. It is not a universal invariant of
cardiac phase synchronization.

\subsection{Early warning signals}

At bifurcation points in the Kuramoto model, the statistical moments of
the order parameter change predictably. We compute four EWS metrics
over a sliding window: variance, lag-1 autocorrelation, skewness, and
kurtosis. These metrics are displayed as a composite risk strip below
the main CSI plot, flagging potential transitions before they manifest
as overt arrhythmias. The EWS metrics are exploratory; their prognostic
value has not been prospectively validated.

\subsection{Respiratory sinus arrhythmia and vagal coherence}

We extract a respiration proxy from ECG baseline wander in the
0.1--0.4 Hz band, then compute coherence between RSA and the
high-frequency component of the RR spectrum. This is the vagal-cardiac
coherence index. It is a proxy for the brain-heart axis derived from
single-lead ECG alone.

\subsection{Apnea screening via coherence band tracking}

Following the reported literature on heart-respiration coherence in
sleep apnea, we track the coherence peak in the 0.2--0.3 Hz band over
time. A vanishing peak is associated with apneic episodes. We report a
per-window apnea likelihood score. This is explicitly a screening aid,
not a diagnosis.

\subsection{Signal quality and runtime verification}

Every metric in this section is only as trustworthy as the signal from
which it was derived. NeuroCardia therefore computes a signal quality
assessment on the raw input, at load time, before any detection or
analysis runs. Four independent sub-metrics are computed:

\begin{enumerate}
\item Signal-to-noise ratio (SNR). The ratio of energy in the QRS band
(5--15 Hz) to energy in the high-frequency noise band (40--100 Hz).
Reported in decibels.
\item Flatline fraction. The proportion of successive sample pairs whose
absolute difference is below an epsilon.
\item Clipping fraction. The proportion of samples at the observed
extremes of the recording.
\item Baseline wander fraction. The ratio of energy below 0.5 Hz to
total energy.
\end{enumerate}

These four are combined into a composite score in $[0, 1]$. Hard floors
apply: if the flatline fraction exceeds 0.9, or the baseline wander
fraction exceeds 0.7, or the SNR is below 0 dB, the composite score is
capped regardless of the other terms. The score maps to a verdict:
GOOD $\geq 0.75 >$ FAIR $\geq 0.50 >$ POOR. The verdict is displayed
in the analyzer screen and included in every report.

This is a runtime verification mechanism, not a validation claim. It
asserts that the input from which the metric was derived meets a
minimum quality threshold, not that the metric is correct.

\section{Validation Strategy}\label{sec:validation}

\subsection{Reference datasets}

To establish analytical validity before clinical trials, NeuroCardia is
tested against three gold-standard annotated corpora from PhysioNet:

\begin{itemize}
\item \textbf{MIT-BIH Arrhythmia Database}~\cite{mitdb}. 48 half-hour
records of ambulatory ECG, annotated beat-by-beat by cardiologists.
\item \textbf{Fantasia Database}. 40 young and elderly healthy subjects
with rigorously screened 2-hour records.
\item \textbf{European ST-T Database}. 90 two-hour recordings focused on
ST and T-wave changes.
\end{itemize}

\subsection{Validation pipeline}

Validation runs as an automated test harness in the CI/CD pipeline. For
every commit: records from the benchmark corpora are streamed through
the app's Kotlin data ingestion pipeline; R-peak annotations produced
by the Pan-Tompkins module are compared against reference
\texttt{.atr} annotation files using a 150 ms match window; the
Kuramoto engine computes $r(t)$ for each record; and summary statistics
are exported to CSV and verified against reference golden outputs.

\subsection{Acceptance criteria}

Each validation test has an explicit numeric target. The targets below
are the commitments made in this white paper.

\begin{table}[h]
\centering
\small
\begin{tabular}{lll}
\toprule
Metric & Dataset & Target \\
\midrule
R-peak sensitivity & MIT-BIH & $\geq 99.0\%$ \\
R-peak PPV & MIT-BIH & $\geq 99.0\%$ \\
HRV (RMSSD) agreement & Fantasia & within $5\%$ \\
HRV (SDNN) agreement & Fantasia & within $5\%$ \\
CSI discriminative validity ($K = 2$) & NSR/AF pilot & AUC $\geq 0.85$ \\
Residual orthogonality & NSR/AF pilot & AUC $\geq 0.65$ after HRV removal \\
EWS prognostic validity & MIT-BIH AF transitions & AUC $\geq 0.70$ \\
Signal quality screen & All datasets & score $\geq 0.5$ \\
\bottomrule
\end{tabular}
\caption{Acceptance criteria. All targets are met by the current
implementation with the exception of the EWS prognostic criterion,
which is exploratory and has not yet been tested on a corpus with
annotated transitions.}
\end{table}

The residual orthogonality criterion is new in revision 0.2. It asks
whether the Kuramoto residual retains diagnostic value after
regressing out the two standard HRV features it is most correlated
with. The current implementation meets the criterion (AUC 0.716 at
$n = 18$), but the margin is narrow and the sample size small.

\subsection{Empirical results}

The CSI depends on two parameters that are not determined by cardiac
physiology: the frequency distribution of the oscillator network, and
the coupling strength $K$. To characterize this dependence and to test
whether the CSI carries information beyond conventional HRV, we ran a
series of experiments against two reference corpora: the MIT-BIH
Arrhythmia Database (48 records, mixed rhythms) for the parameter
sensitivity analysis, and a paired sample of MIT-BIH Normal Sinus
Rhythm (9 records) and MIT-BIH Atrial Fibrillation (9 records) for the
discriminative validity test.

\subsubsection{Sensitivity to the frequency distribution}

Five frequency distributions were tested: Riemann zeros (Arm A, the
production configuration), uniform random (Arm B), Gaussian (Arm C),
physiological scale (Arm D), and linear spacing (Arm E). Table~\ref{tab:trajectory-wp}
reports the mean Pearson correlation between Arm A's trajectory and
each alternative arm's trajectory.

\begin{table}[h]
\centering
\begin{tabular}{r|cccc}
\toprule
$K$ & Arm B & Arm C & Arm D & Arm E \\
\midrule
2.0 & $0.052$ & $0.076$ & $0.051$ & $0.168$ \\
5.0 & $0.250$ & $0.230$ & $0.250$ & $0.611$ \\
10.0 & $0.386$ & $0.366$ & $0.386$ & $0.778$ \\
30.0 & $0.458$ & $0.428$ & $0.458$ & $\mathbf{0.869}$ \\
100.0 & $-0.000$ & $0.000$ & $0.001$ & $0.001$ \\
\bottomrule
\end{tabular}
\caption{Mean Pearson correlation between the CSI trajectory under Arm
A and under each alternative arm, across 48 mitdb records. The
pre-registered threshold of $0.80$ is met only by Arm E at $K = 30$,
where the network is saturated.}
\label{tab:trajectory-wp}
\end{table}

The CSI trajectory is not robust to the frequency distribution at any
coupling strength where the metric has dynamic range. The companion
sensitivity analysis reports the full result.

\subsubsection{Sensitivity to coupling strength}

The coupling strength was swept over $K \in \{0.5, 1, 2, 3, 4, 5, 10\}$.
The residual diagnostic peaks at $K = 2$ (AUC 0.89), inverts below
$K = 1$ (AUC 0.17--0.21), and decays to noise above $K = 3$
(AUC 0.59--0.72). The useful window is narrow.

\subsubsection{Correlation with conventional HRV}

At $K = 10$, the mean CSI correlates with the coefficient of variation
of the RR interval series at $\rho = -0.95$, with SDNN at
$\rho = -0.87$, and with the Poincar\'e SD2 at $\rho = -0.86$. The CSI
is essentially a monotone function of RR variability at this
parameterization.

At $K = 2$, the CSI is not a monotone function of RR variability, and
the residual (the component orthogonal to RR variability) separates
sinus from AF.

\subsubsection{Residual analysis}

We regressed CSI mean against RR coefficient of variation by ordinary
least squares and analyzed the residuals. Table~\ref{tab:residual-wp}
reports the residual AUC at $K = 2$ for each of the five frequency
distributions.

\begin{table}[h]
\centering
\begin{tabular}{llrr}
\toprule
Arm & Distribution & Residual AUC & $p$ \\
\midrule
A & Riemann zeros & $0.889$ & $0.002$ \\
B & Uniform random & $0.543$ & $0.986$ \\
C & Gaussian & $0.852$ & $0.007$ \\
D & Physiological scale & $0.556$ & $0.982$ \\
E & Linear spacing & $0.840$ & $0.010$ \\
\bottomrule
\end{tabular}
\caption{Residual AUC at $K = 2$ for each frequency distribution, on
the paired sinus/AF pilot corpus. Three arms exhibit the residual; two
do not.}
\label{tab:residual-wp}
\end{table}

The residual is diagnostic for three of five tested distributions.
Revision 0.1 of this white paper reported only the three working arms
and described the result as robust across the frequency distribution;
the full five-arm test shows that this claim is narrower than
reported.

\subsubsection{Orthogonality of the residual to standard HRV}

To test whether the residual carries information beyond standard HRV,
we regressed it on the two features it is most correlated with in the
pilot corpus (the Poincar\'e SD1/SD2 ratio and the coefficient of
variation of the RR series) and computed the AUC of the
orthogonalized residual. The raw residual separates sinus from AF with
AUC 0.889; the orthogonalized residual separates them with AUC 0.716.
Approximately two thirds of the residual's diagnostic value is
attributable to its linear relationship with the two HRV features; one
third is not.

At $n = 18$ the confidence interval on the orthogonalized AUC is wide
and the difference between 0.889 and 0.716 is not statistically
significant at this sample size. The finding is suggestive, not
conclusive.

\subsubsection{Mechanism}

In the AF cohort, 94\% of the residual's \emph{between-record} variance
is explained by three P-wave features extracted from the raw ECG
(P-wave energy, PR-interval variability, and fibrillatory amplitude),
features the Kuramoto engine does not have access to. In the sinus
cohort, 37\% is explained. The residual is best understood as a
between-record discriminator, not a within-record state tracker: the
window-to-window variation of the residual within a single record does
not correlate with the same features (mean $|r| < 0.15$ in both
cohorts), and the window-level AUC plateaus at 0.66--0.68 across
1--5 minute windows (0.50 at 30~s). Whole-record averaging is required
to recover the discriminative signal. Two thirds of the residual's
record-level diagnostic value is redundant with standard HRV.

\subsubsection{Interpretation}

Three conclusions follow.

First, the CSI at $K = 2$ carries information about cardiac rhythm that
is not captured by conventional HRV. After regressing out the linear
dependence on RR variability, the residual separates normal sinus
rhythm from atrial fibrillation with an AUC of 0.89 on the pilot
corpus. After regressing out both RR variability and the Poincar\'e
SD1/SD2 ratio, the residual retains an AUC of 0.716.

Second, the metric's behavior depends critically on the coupling
strength. The same metric at $K = 10$ is redundant with conventional
HRV, with no residual diagnostic information. Above $K = 3$ the network
saturates; below $K = 1$ the sign of the residual correlation inverts.
The useful regime is narrow.

Third, the residual is not distribution-agnostic. It is exhibited by
three of five tested frequency distributions. The choice of distribution
matters.

\subsection{Reproducibility}

All validation code and expected outputs are committed to the
repository under \texttt{validation/}. Every metric is reproducible
with a single Gradle command on any JVM:
\begin{verbatim}
./gradlew :validation:run --args="<mode> <args>"
\end{verbatim}
The output is deterministic given the same input and the same build.
Builds are signed, and the signed APK hash is published with every
release.

\section{Regulatory Pathway (FDA 510(k))}\label{sec:regulatory}

\subsection{Intended use and device classification}

NeuroCardia is intended to be used as a software-as-a-medical-device
(SaMD) decision-support tool. It processes single-lead ECG data to
compute heart rate variability metrics and phase-domain synchronization
indices. It is intended for use by healthcare professionals and trained
investigators in research and clinical settings. It does not provide an
automated diagnosis, does not direct clinical treatment, and is not
intended as a replacement for standard ECG interpretation.

Under FDA product classification, ECG analysis software of this type
falls under product code DQK (electrocardiograph software) or MWJ
(system, network, and patient monitoring), typically regulated as
Class II medical devices requiring 510(k) clearance.

\subsection{Predicate device selection}

The primary predicate device for 510(k) clearance will be an
established commercial ECG analysis software package with cleared
indications for automated HRV analysis and arrhythmia detection from
ambulatory ECG recordings.

\subsection{Substantial equivalence argument}

NeuroCardia is substantially equivalent in intended use and fundamental
technological characteristics to predicate devices. The technological
difference --- the introduction of the Kuramoto coupled-oscillator
model and the derived CSI $r(t)$ --- represents a new feature for a
cleared SaMD. Under FDA guidance, new analytical methods in an existing
device category are addressed through analytical validation and
non-clinical performance testing, without necessarily requiring a
De Novo pathway, provided the device does not make novel diagnostic
claims that alter its risk profile.

\subsection{FDA classification strategy}

We plan a 510(k) premarket notification pathway. The claims made in the
initial submission are: R-peak detection and heart rate computation;
HRV metrics (time and frequency domain); rhythm irregularity flagging;
and report export (PDF) suitable for clinical documentation.

Claims that require additional validation before regulatory submission:
cardiac synchronization index as a diagnostic; early warning signal
detection; apnea screening; vagal-cardiac coherence index; and residual
orthogonality to standard HRV.

These are initially released as research features, gated behind a
settings flag, and documented as not-for-clinical-use. A subsequent
510(k) or De Novo submission will cover them once validation is
complete.

\subsection{Software lifecycle and quality management}

To comply with FDA expectations for SaMD (and IEC 62304 standard for
medical device software lifecycle processes), development adheres to
version control and reproducibility (every release is tagged,
cryptographically hashed, and tied to a Software Bill of Materials in
CycloneDX format); traceability (requirements are mapped to unit tests
and integration tests); risk management (ISO 14971; identified hazards
such as false R-peak classification are mitigated through signal
quality checks, confidence scoring, and mandatory clinician review
disclaimers); and quality management (ISO 13485-aligned QMS before any
clinical claim).

\section{Limitations and Open Questions}\label{sec:limitations}

\subsection{The Riemann-zero frequency distribution}

The author does not claim that the Riemann zeros have a direct
physiological correlate in cardiac electrophysiology. The choice is a
specific choice documented in the companion mathematical reference.
Whether this distribution offers diagnostic advantages over standard
Gaussian or uniform frequency distributions is addressed empirically in
Section~5.4. The sensitivity analysis shows that the CSI trajectory is
not invariant to the choice of frequency distribution at any tested
coupling strength, and that the residual diagnostic is exhibited by
three of five tested distributions.

\subsection{Carrier scale and phase velocity}

The reference parameterization places the carrier frequency at
$\Omega_T = 587{,}000$ rad/s, five orders of magnitude above the cardiac
rhythm. The absolute scale of the carrier does not affect the
dimensionless order parameter $r(t)$, but it has a decisive effect on
numerical integration. We work with normalized frequencies, in which
the carrier is subtracted and the residual Riemann-zero distribution is
scaled to unit variance.

\subsection{Signal quality and artifact sensitivity}

Like all HRV and phase-domain tools, NeuroCardia is sensitive to motion
artifact, electrode displacement, and ectopic beats. While the
Pan-Tompkins detector includes amplitude and slope filtering, severe
baseline wander or frequent premature ventricular contractions can
distort the RR tachogram.

\subsection{Choice of $N = 16$}

The number of oscillators is a free parameter. We adopt $N = 16$ to
match the reference implementation. The sensitivity of the CSI to $N$
is a planned study.

\subsection{Single-lead constraints}

The absence of a spatial map of the atrium limits rotor detection to a
temporal analysis. A multi-lead device would allow spatial mapping.
That is a different device with a different intended use.

\subsection{Residual orthogonality}

The orthogonality result (AUC 0.716 at $n = 18$) has a wide confidence
interval. Full-corpus replication on the complete NSRDB and AFDB is
required before the residual can be claimed as a validated
discriminative feature.

\subsection{Validation is ongoing}

At the time of writing, the novel metrics are supported by literature
and by internal consistency checks, but not by prospective clinical
validation. This document will be revised as validation results become
available. Readers should treat every claim in Section~\ref{sec:metrics}
as a hypothesis under test, not as an established result.

\section{Roadmap and Future Work}\label{sec:roadmap}

Development milestones for NeuroCardia are structured across five
version releases:

\begin{itemize}
\item \textbf{v0.1 (Current).} Core Android app, Pan-Tompkins R-peak
detector, basic time-domain HRV, Kuramoto engine with $N = 16$ and
Riemann zero frequencies, CSV (RR intervals) and PDF report export,
GPL-3.0 open-source release.
\item \textbf{v0.2.} EWS metrics (variance, autocorrelation), CSV
exports for metrics and annotated signals, JSON audit log and SHA-256
manifest generation.
\item \textbf{v0.3.} Single-file HTML reports, session bundle JSON
schema (v1), automated SBOM generation.
\item \textbf{v0.4.} Full $r(t)$ CSI trajectory export, advanced
coupling topology selectors, companion UI plots.
\item \textbf{v0.5.} WFDB (.dat + .hea + .atr) and EDF+ export support
for direct interoperability with PhysioNet and clinical workstations.
\end{itemize}

\section{Conclusion}

NeuroCardia is an open-source Android application for offline,
on-device analysis of single-lead ECG recordings. It runs entirely on
the handset, with no network dependency, no user accounts, and no
telemetry. It computes conventional HRV metrics, a four-axis signal
quality score, and a phase-domain Cardiac Synchronization Index derived
from a Kuramoto coupled-oscillator network. Its DSP and I/O layers are
packaged as standalone Kotlin/JVM modules that can be reused in any JVM
project.

The application's primary contribution is the deployment envelope:
open-source, offline, on-device ECG analysis on commodity Android
hardware, with a regulatory strategy that treats the conventional HRV
metrics as validated outputs and the phase-domain metric as a research
feature with a documented plan for further validation. The
phase-domain metric itself is a record-level nonlinear summary of the
RR series with partial novelty over standard HRV (approximately one
third of its diagnostic value is not explained by the two HRV features
it is most correlated with) and with specific, quantified limitations
(three of five tested frequency distributions; record-level but not
window-level discrimination). These properties are documented in the
companion residual paper.

The version 0.1 white paper is revised here as version 0.2. The
principal changes are a reframing of the application's contribution
around the deployment envelope rather than the novel metric, the
addition of an explicit "what this application is not" section, and
the scoping of the phase-domain claims to their empirically
demonstrated range. The regulatory strategy is unchanged.

\begin{thebibliography}{9}

\bibitem{taskforce1996}
Task Force of the European Society of Cardiology and the North American
Society of Pacing and Electrophysiology. Heart rate variability:
standards of measurement, physiological interpretation, and clinical
use. \emph{Circulation}, 93(5):1043--1065, 1996.

\bibitem{kuramoto1975}
Y.~Kuramoto. Self-entrainment of a population of coupled non-linear
oscillators. In \emph{International Symposium on Mathematical Problems
in Theoretical Physics}, pages 420--422. Springer, 1975.

\bibitem{acebron2005}
J.~A.~Acebr\'on, L.~L.~Bonilla, C.~J.~P\'erez Vicente, F.~Ritort, and
R.~Spigler. The Kuramoto model: A simple paradigm for synchronization
phenomena. \emph{Reviews of Modern Physics}, 77(1):137--185, 2005.

\bibitem{pan1985}
J.~Pan and W.~J.~Tompkins. A real-time QRS detection algorithm.
\emph{IEEE Transactions on Biomedical Engineering},
BME-32(3):230--236, 1985.

\bibitem{mitdb}
G.~B.~Moody and R.~G.~Mark. The impact of the MIT-BIH Arrhythmia
Database. \emph{IEEE Engineering in Medicine and Biology Magazine},
20(3):45--50, 2001.

\bibitem{nsrdb}
R.~G.~Mark and G.~B.~Moody. MIT-BIH Normal Sinus Rhythm Database.
PhysioNet, 1990. \url{https://physionet.org/content/nsrdb/}.

\bibitem{afdb}
G.~B.~Moody and R.~G.~Mark. MIT-BIH Atrial Fibrillation Database.
PhysioNet, 1983. \url{https://physionet.org/content/afdb/}.

\bibitem{scheffer2009}
M.~Scheffer et al. Early-warning signals for critical transitions.
\emph{Nature}, 461(7260):53--59, 2009.

\bibitem{simoes2026engine}
A.~Simoes. The Kuramoto Engine: Physics and Mathematics of the Cardiac
Synchronization Index. Technical reference, revision 0.2, 2026.
Companion document.

\bibitem{simoes2026sensitivity}
A.~Simoes. Sensitivity Analysis Protocol: Validation Methodology for
the Kuramoto Phase Engine in NeuroCardia. Protocol, revision 0.2,
2026. Companion document.

\bibitem{simoes2026residual}
A.~Simoes. The Kuramoto Residual: Phase-Domain Information in the
Cardiac Rhythm Beyond Heart Rate Variability. Research paper, revision
0.2, 2026. Companion document.

\end{thebibliography}

\end{document}
